\section{Generative Models in Biomedical Contexts}
\label{sec:methods}

The application of deep generative models to biomedical data synthesis requires careful consideration of architectural choices, training dynamics, and domain-specific constraints. This section examines the technical foundations of three predominant approaches---Generative Adversarial Networks (GANs), Variational Autoencoders (VAEs), and diffusion models---with emphasis on adaptations that address the unique characteristics of rare disease datasets.

\subsection{Generative Adversarial Networks}

GANs employ adversarial training between two neural networks: a generator $G$ that produces synthetic samples from random noise, and a discriminator $D$ that distinguishes real from generated data. The training objective formulates a minimax game where the generator seeks to minimize the discriminator's classification accuracy while the discriminator maximizes it. This adversarial process, when successful, drives the generator to produce samples indistinguishable from the real data distribution.

Frid-Adar et al. demonstrated the application of Deep Convolutional GANs (DCGANs) to medical imaging, implementing a generator architecture with four fractionally-strided convolutional layers using 5×5 kernels \cite{fridadar2018gan}. The network transforms 100-dimensional uniform noise vectors into 64×64 pixel liver lesion images through progressive upsampling. Batch normalization and ReLU activations stabilize training, while the final layer employs tanh activation to constrain output values to the appropriate intensity range. The discriminator mirrors this architecture in reverse, classifying input images as real or synthetic through successive downsampling operations.

A critical challenge in GAN training is mode collapse, where the generator produces limited diversity by concentrating on a subset of the target distribution. Choi et al. introduced minibatch averaging as an effective mitigation strategy for electronic health record synthesis \cite{choi2017generating}. Rather than processing individual samples independently, the discriminator computes feature statistics across minibatches and incorporates this aggregated information into its classification decision. This technique encourages the generator to produce diverse outputs without the hyperparameter sensitivity of alternative approaches such as minibatch discrimination. The medGAN architecture further incorporates shortcut connections between generator layers, enabling more effective gradient flow during backpropagation and accelerating convergence.

The choice between different GAN variants depends on the synthesis task requirements. Auxiliary Classifier GANs (ACGANs) extend the discriminator to predict class labels in addition to authenticity, enabling conditional generation of specific disease subtypes. Frid-Adar et al. found that DCGAN outperformed ACGAN for their liver lesion classification task, suggesting that the additional auxiliary loss may not always improve performance when category information is already well-represented in the training data \cite{fridadar2018gan}. Conditional GANs (cGANs) and their variants---including CycleGANs for unpaired image translation and StyleGANs for fine-grained control over generated features---have each found distinct applications in biomedical synthesis \cite{mendes2025synthetic}.

\subsection{Variational Autoencoders}

VAEs approach generative modeling through probabilistic inference, learning a latent representation that captures the underlying data distribution. The architecture comprises an encoder network that maps input data to a probability distribution over a lower-dimensional latent space, and a decoder that reconstructs samples from latent vectors. Unlike GANs' adversarial objective, VAEs optimize a variational lower bound on the data log-likelihood, combining a reconstruction loss with a Kullback-Leibler divergence term that regularizes the latent distribution toward a standard normal prior.

This probabilistic formulation provides theoretical guarantees on the quality of the learned latent space and eliminates the training instability often encountered with GANs. The regularization term ensures that similar data points map to nearby latent representations, facilitating interpolation between samples and enabling controlled generation through latent space navigation. Conditional VAEs (CVAEs) extend this framework by conditioning both encoder and decoder on auxiliary information such as disease labels, demonstrating particular effectiveness with smaller datasets that require class-specific generation \cite{mendes2025synthetic}.

The primary limitation of VAEs relative to GANs lies in reconstruction quality. The continuous latent space and probabilistic decoder produce samples that, while statistically faithful to the training distribution, often exhibit reduced sharpness compared to adversarially-trained generators. For medical imaging applications requiring high spatial resolution, this trade-off manifests as slightly blurred synthetic images. However, VAEs require substantially lower computational resources than GANs, making them viable for resource-constrained environments \cite{mendes2025synthetic}. Hybrid VAE-GAN architectures attempt to combine the advantages of both approaches, using VAE structure for stable training while incorporating adversarial objectives to sharpen generated outputs.

\subsection{Diffusion Models}

Diffusion probabilistic models represent an emerging paradigm for generative modeling, formulating synthesis as a progressive denoising process. Training proceeds in two phases: a forward diffusion process gradually adds Gaussian noise to real data over $T$ timesteps until reaching pure noise, followed by learning a reverse process that reconstructs clean data from noisy inputs. The reverse process, parameterized by a neural network, predicts the noise component at each timestep, enabling iterative refinement from random samples to high-quality outputs.

Zhong et al. adapted this framework to sequential electronic health record data through the MedDiffusion architecture \cite{zhong2024meddiffusion}. The model embeds medical diagnosis codes and temporal information into visit representations using a combination of learned embeddings and gap-based time encodings. A bidirectional Long Short-Term Memory (LSTM) network processes these visit sequences to capture longitudinal patient trajectories. The key innovation lies in the step-wise attention mechanism, which aggregates information from the current visit and previous hidden states to condition the diffusion process on patient history. This conditioning enables generation of synthetic visits that maintain realistic disease progression patterns while expanding the training dataset.

The forward diffusion process in MedDiffusion adds noise to visit embeddings according to a predetermined variance schedule, gradually corrupting the structured medical code representations. The reverse process trains a denoising network to predict the original clean embedding from noisy inputs at each timestep, conditioned on the LSTM-encoded patient history. Generating a new synthetic visit begins with Gaussian noise and applies the learned denoising steps iteratively, with the LSTM context ensuring temporal consistency across the patient trajectory.

Zhong et al. demonstrated that this diffusion-based approach outperformed GAN baselines (ehrGAN, TabDDPM) across multiple health risk prediction tasks, achieving 77.88\% PR-AUC on kidney disease prediction compared to 74.92\% for the best GAN baseline \cite{zhong2024meddiffusion}. The superior performance stems from diffusion models' ability to generate in continuous latent space while avoiding the mode collapse and training instability that plague GAN training on sequential discrete data. Ablation studies confirmed that both the step-wise attention mechanism and the diffusion-based generation contribute essential components to the overall performance.

\subsection{Comparative Analysis}

The choice of generative model architecture depends on multiple factors including data modality, available computational resources, dataset size, and synthesis objectives. Medical imaging applications, which dominated 70.3\% of reviewed studies, predominantly employ GANs due to their capacity for generating sharp, high-resolution images \cite{finetti2025data}. The adversarial training objective proves particularly effective for capturing the complex spatial dependencies present in CT scans, MRI sequences, and histopathology slides. However, this comes at substantial computational cost---GAN training typically requires GPU clusters and careful hyperparameter tuning to achieve stable convergence.

VAEs offer a compelling alternative when researchers face computational resource constraints or require theoretical guarantees on the latent space structure. Their lower training complexity and reduced risk of mode collapse make them particularly suitable for exploratory analysis and scenarios where slightly reduced image sharpness is acceptable relative to training efficiency. The probabilistic framework also facilitates uncertainty quantification, enabling assessment of the model's confidence in generated samples.

Diffusion models have emerged most recently and show particular promise for sequential and temporal data where maintaining consistency across time steps is critical. The iterative denoising process naturally incorporates contextual information, making diffusion models well-suited to electronic health record synthesis and longitudinal patient trajectory generation. While diffusion models require more inference steps than single-pass GAN or VAE generation, recent advances in accelerated sampling have reduced this computational overhead substantially \cite{zhong2024meddiffusion}.

Hybrid approaches combining elements from multiple frameworks represent an active area of investigation. VAE-GAN architectures leverage VAE encoders for stable latent space learning while incorporating GAN discriminators to sharpen outputs. Conditional variants of all three model classes enable class-specific generation, proving essential for rare disease applications where maintaining disease subtype distinctions is critical. The continued rapid development of generative modeling techniques suggests that optimal architectures for biomedical synthesis remain an evolving target rather than a settled question.
